\subsection*{Paso 1. Entender qué aspecto tiene la infraestructura moderna}

El estándar de la industria para infraestructuras de red empresariales se fundamenta en un diseño jerárquico. Este modelo, que comúnmente incluye las capas de \textit{Acceso}, \textit{Distribución} y \textit{Núcleo} (o un núcleo colapsado), está diseñado para garantizar escalabilidad, modularidad y predictibilidad en el rendimiento del tráfico al definir funciones claras para cada estrato de la red.

Una característica fundamental de las redes modernas es la evasión de topologías ``planas''. El uso de Redes de Área Local Virtuales (\texttt{VLANs}) es el mecanismo estándar para la segmentación de dominios de broadcast. Esta segmentación es crucial para aislar lógicamente el tráfico por función (ej. datos de usuario, voz sobre IP, gestión de red, tráfico de invitados o videovigilancia) y sirve como base indispensable para la aplicación de políticas de seguridad granulares mediante enrutamiento inter-VLAN y Listas de Control de Acceso (\texttt{ACLs}).

En términos de rendimiento y disponibilidad, la infraestructura moderna debe soportar conectividad inalámbrica de alta densidad, empleando estándares actuales como \texttt{Wi-Fi 6 (802.11ax)} o \texttt{Wi-Fi 7 (802.11be)}, para manejar la proliferación de dispositivos móviles. Asimismo, la Alta Disponibilidad (HA) es un requisito, implementada a través de mecanismos de redundancia en enlaces físicos (Agregación de Enlaces - \texttt{LACP}) y protocolos de conmutación por error para asegurar la continuidad del servicio.

Finalmente, la gestionabilidad es un pilar esencial. La infraestructura debe ser observable (monitoreable) a través de protocolos estándar como \texttt{SNMP (Simple Network Management Protocol)} y análisis de flujo (como \texttt{NetFlow}), y ser susceptible a una gestión centralizada. Esta capacidad de gestión es el cimiento sobre el cual se construyen las arquitecturas de automatización y Redes Definidas por Software (\texttt{SDN}).